\section{终局取舍：效率是表象，状态接口才是本质}

前面的模型、token 与 H-Net 例子都服务于最终修正：SSM 不只是“更快但更弱”，attention 也不只是“更强但二次”。两者的优缺点由 autoregressive state 的形式共同产生。最后四张图先重申 attention 的适用分辨率，再分别改写 SSM 与 attention 的 tradeoff，最终把架构研究压缩成 FLOPs-to-capabilities 问题。

\subsection{Attention 真正强在 pre-compressed abstraction 上}

本节回到整堂课反复出现的 abstraction 问题。漫画把 Transformer 在 raw high-resolution input 上的尴尬画成误会：模型并非天然理解任何粒度，它在有意义、可组合、可单独寻址的单位上最强。图右侧红字不是否定 attention，而是在给它定位——先把数据压到合适层，再利用全局内容寻址。

\lecturefigure{slide-030.jpg}{Attention 最有效的区域：pre-compressed data 与正确 abstraction level}{Stanford 官方录像，00:49:20--00:50:09}

\begin{importantbox}{架构选择应该是 resolution-aware 的}
同一模态可以在不同层使用不同 primitive：像素/byte 层先做 local or recurrent compression，object/word/chunk 层再做 attention，最终决策层可能再接 external memory。把整个 pipeline 强制成单一 layer type，往往是在让架构适应错误的单位。
\end{importantbox}

\subsection{SSM：把 efficiency 划掉，保留 statefulness 与 compression}

图中先写“linear-time scaling, constant-time inference”，随后把 efficiency 划掉，换成 stateful / compressive model，并标注 brain-like state。讲者不是说效率不重要，而是说它不足以解释为什么这种模型在 state tracking、online interaction 与 abstraction building 上可能更合适；这些能力与 retrieval weakness 都来自有限状态。

\lecturefigure{slide-031.jpg}{最终 SSM 取舍：stateful compression、state tracking 与 fine-grained retrieval loss}{Stanford 官方录像，00:50:10--00:51:19}

\teachervoice{00:50:05--00:51:27，Gu 说 efficiency 是 red herring：SSM 的 pros and cons 是同一枚硬币的两面，即 brain-like autoregressive state。有限状态既丢精确细节，也创造追踪与抽象的压力。}

\subsection{Attention：把 quadratic 划掉，保留 resolution dependence}

Attention 的强项是对过去 context 做 fine-grained access，尤其适合 recall 与 retrieval。图中把“quadratic scaling”划掉，不是取消复杂度，而是指出更深的限制：模型对输入 token 的分辨率与语义高度敏感，database-like cache 会忠实保存它被给到的单位，却不会替设计者决定这些单位是否值得保存。

\lecturefigure{slide-032.jpg}{最终 attention 取舍：database-like cache、精确检索与 token-resolution dependence}{Stanford 官方录像，00:51:20--00:53:09}

\begin{warningbox}{Quadratic attention 在需要精确记忆时并不“浪费”}
若任务确实要求逐项保留合同条款、代码 token、证据段落或外部工具返回，$O(L)$ cache 与 $O(L)$ 每步读取可能是完成任务所需的信息成本。正确问题不是“能否消灭 quadratic”，而是“哪些层、哪些 token、哪些时间尺度值得支付这个成本”。
\end{warningbox}

\begin{table}[H]
\centering
\caption{两种 autoregressive state 的对称取舍}
\begin{tabularx}{\textwidth}{L{0.18\textwidth}>{\raggedright\arraybackslash}X >{\raggedright\arraybackslash}X}
\toprule
比较轴 & Token-resolution attention state & Compressed recurrent state \\
\midrule
跨步状态 & 随 $L$ 增长的 KV cache & 与 $L$ 解耦的有限 state \\
精确访问 & 强；历史位置可单独寻址 & 弱；细节混入 shared state \\
在线更新 & 每步读取历史 cache & 固定接口递推更新 \\
归纳偏置 & database、retrieval、copy & tracking、smoothing、compression \\
敏感因素 & token resolution 与 semantic unit & state capacity 与 update expressivity \\
典型补救 & MQA/GQA、KV compression、retrieval routing & 扩大 state、selectivity、插入 attention \\
\bottomrule
\end{tabularx}
\end{table}

\subsection{架构研究的中心问题：每个 FLOP 是否花在有意义的结构上}

最后一张图把训练系统画成黑箱：左边投入 FLOPs 与数据，右边得到 capabilities。真正目标不是最小 FLOPs，也不是最大模型，而是提高 compute 到能力的转换效率。若 Transformer 在 character stream 中缓存每个字符却很少需要单字符寻址，这些 FLOPs 没有匹配任务结构；若任务必须精确回忆每个细节，支付 attention 成本反而合理。

\lecturefigure{slide-033.jpg}{FLOPs $\rightarrow$ model $\rightarrow$ capabilities：模型是否明智地使用计算？}{Stanford 官方录像，00:53:10--00:54:34}

\[
\text{system efficiency}
=\frac{\text{useful capabilities}}{\text{training dollars}}
=\frac{\text{FLOPs}}{\text{dollars}}
\times
\frac{\text{capabilities}}{\text{FLOPs}}.
\]
第一项受 hardware、kernel、parallelism 与 utilization 影响；第二项受 architecture、data、objective 与 representation 影响。Gu 主要研究第二项，但部署系统必须同时优化两项。一个 FLOP 在数学上相同，在信息结构上却可能完全不同：它可以用于比较有意义 chunk，也可以用于反复处理无意义高分辨率单位。

\begin{importantbox}{最终设计准则}
不要先选架构再强迫数据适配，而应先画出信息在不同时间尺度上的结构：哪些细节必须保留，哪些可压缩，哪些需要随机访问，哪些只需状态追踪。然后把 attention、SSM、convolution、local encoder 与 external memory 放到最匹配的层级。
\end{importantbox}

\subsection{Q\&A：backprop、硬件与“局部最优”}

一位同学质疑：现代 SSM 为了 backprop-through-time 与 GPU 并行牺牲 recurrence expressivity，而人脑并不保存许多副本做反向传播。Gu 的回答很谨慎：他认为当前架构、backprop 与 hardware 可能处在共同 local optimum；backprop 的 activation memory 会对可学习 dependency length 施加物理限制，但现实任务是否已能通过工程手段绕过，仍不确定。

\teachervoice{00:55:10--00:57:22，Gu 认为重新设想 learning algorithm 与 hardware constraint 可能产生 fundamentally better models；同时承认自己没有给出可替代 backprop 的成熟方案。}

\begin{warningbox}{开放问题不等于已证实缺陷}
“人脑不用 backprop”不能推出 backprop 对机器学习无效；“显存装不下超长 sequence”也不等于模型绝对学不到长依赖，因为 curriculum、synthetic task、retrieval 与 recurrent training 都可能提供间接路径。课堂观点应作为研究假设，而不是结论。
\end{warningbox}

\subsection{Q\&A：H-Net 的边界、memory 与 tiny models}

关于 H-Net，讲者明确说当前版本只是 first working step，尚未在最大尺度验证。动态 chunk 没有无限 vocabulary：初始 byte 进入有限 embedding，后续全部在 latent space 中选择 boundary、汇聚 chunk representation。对于 long context，他设想递归压缩成越来越高层的 memory，再让 inner model attention 更少、更有意义的条目，但尚无后续验证。

\teachervoice{00:59:09--01:03:31，chunk representation 不通过词表 lookup，而是在 encoder latent space 中保留被选中的位置或 pooled summary；hierarchical memory 是与 H-Net philosophy 一致的未来方向，不是当前结果。}

对 data-constrained / tiny model，Gu 直接回答“不知道一般规律”。他只提到当时一个很小的 Mamba-3 experiment 仍发现少量 attention 很关键。这种不确定性本身值得保留：归纳偏置在不同 scaling regime 可能改变，不能从 billion-parameter curve 反推 tiny model。

\subsection{Q\&A：SSM 能否像 attention map 一样解释}

某些 SSM / linear-attention 形式可展开为 token-to-token dependence map，视觉上类似 attention map。Gu 的观察是，softmax attention 有时表现为近似 hard attention，只集中到少数 token；SSM dependence 往往更 diffuse，符合“把信息揉进状态”的直觉。但这只是一个 proxy，mechanistic interpretability 仍在进行。

\teachervoice{01:04:52--01:06:29，讲者说 attention-like visualization 可以揭示两类模型的行为差异，却不是唯一解释工具；更完整的 SSM mechanistic interpretability 仍是 open work。}

\subsection{本章小结}

终局比较不再是“二次模型 versus 线性模型”，而是 database-like cache versus brain-like compressed state。Attention 用逐 token 成本换取精确访问，SSM 用不可逆压缩换取在线状态与抽象压力；hybrid 与 hierarchy 让两者按分辨率分工。架构研究最终要优化的，是每个 FLOP 与任务信息结构的匹配程度。

